\documentclass[11pt,a4paper]{article}
\usepackage[T1]{fontenc}
\usepackage[utf8]{inputenc}
\usepackage[english]{babel}
\usepackage{amsmath,amssymb}
\usepackage[margin=2cm]{geometry}
\setlength{\parindent}{0pt}
\setlength{\emergencystretch}{3em}
\allowdisplaybreaks

\title{About \texttt{equivalence\_relations.tex}}
\author{}
\date{}
\begin{document}
\maketitle
This note explains what the test data in \texttt{equivalence\_relations.tex} is used for. That file is input for the engine in \texttt{main.c}; it is read by the test \texttt{test\_maps\_and\_relations.c} in this folder.

\begin{enumerate}
\item \textbf{Purpose.} The file holds 28 statements about equivalence relations and partitions, each to be proved. The statements are data: the engine knows none of them in advance.

\item \textbf{What is being tested.} Whether a given relation is an equivalence relation, as for $a - b \in \mathbb{Z}$ on $\mathbb{Q}$; general facts about equivalence relations; and partitions built from other partitions.

\item \textbf{How the statements are written.} Every property is written out as a formula, not named: for instance ``$f$ is injective'' is $\forall a, b\ (f(a) = f(b) \implies a = b)$. The engine must work from these definitions.

\item \textbf{How to run it.} From the top folder of the repository:
\begin{verbatim}
cc -std=c11 -O2 -Wall -Wextra -o /tmp/test_maps_and_relations \
    tests/set_theory/test_maps_and_relations.c && /tmp/test_maps_and_relations
\end{verbatim}

\item \textbf{Expected result.} One \texttt{PASS} line for the file when every statement is proved with a written proof. The engine does not yet treat statements of this kind: each one is reported as not supported, and the test fails. The file states the goal.
\end{enumerate}
\end{document}
